% Bagian: model-theory
% Bab: interpolation
% Subbagian: definability

\documentclass[../../../include/open-logic-section]{subfiles}

\begin{document}

\olfileid[id]{mod}{int}{def}

\olsection{Teorema Keterdefinisian}

Salah satu penerapan penting teorema interpolasi adalah teorema
keterdefinisian Beth. Untuk mendefinisikan relasi $n$-tempat~$R$, kita dapat
memberikan !!a{formula}~$!C$ dengan $n$ !!{variable} bebas yang tidak
melibatkan~$R$. Ini merupakan definisi \emph{eksplisit} atas~$R$ dalam
kerangka~$!C$. Selanjutnya, kita juga dapat mengatakan bahwa suatu
teori~$\Sigma(P)$ dalam !!{language} yang memuat !!{predicate}
beraritas-$n$~$P$ mendefinisikan~$P$ secara eksplisit jika teori tersebut
memuat (atau setidaknya mengakibatkan) suatu definisi eksplisit yang
diformalkan, yakni
\[
\Sigma(P) \Entails \lforall[x_1][\dots
  \lforall[x_n][(\Atom{P}{x_1,\dots, x_n} \liff !C(x_1, \dots,
    x_n))]].
\]
Namun, definisi eksplisit hanyalah salah satu cara mendefinisikan---dalam arti
menentukan sepenuhnya---suatu relasi. Suatu teori juga dapat bersifat
sedemikian rupa sehingga interpretasi~$P$ ditetapkan oleh interpretasi bagian
lain dari !!{language} dalam setiap model. Teorema keterdefinisian menyatakan
bahwa setiap kali suatu teori menetapkan interpretasi~$P$ dengan cara
ini---setiap kali teori itu \emph{mendefinisikan $P$ secara implisit}---teori
tersebut juga mendefinisikannya secara eksplisit.

\begin{defn}
Andaikan $\Lang{L}$ adalah !!a{language} yang tidak memuat
!!{predicate}~$P$. Suatu himpunan $\Sigma(P)$ dari !!{sentence}s dalam
$\Lang{L} \cup \{P\}$ \emph{mendefinisikan $P$ secara eksplisit} jika dan
hanya jika terdapat !!a{formula}~$!C(x_1, \dots, x_n)$ dari $\Lang{L}$
sedemikian sehingga
\[
\Sigma(P) \Entails \lforall[x_1][\dots
  \lforall[x_n][(\Atom{P}{x_1,\dots, x_n} \liff !C(x_1, \dots,
    x_n))]].
\]
\end{defn}

\begin{defn}
Andaikan $\Lang{L}$ adalah !!a{language} yang tidak memuat
!!{predicate}s~$P$ dan~$P'$. Suatu himpunan $\Sigma(P)$ dari !!{sentence}s
dalam $\Lang{L} \cup \{P\}$ \emph{mendefinisikan $P$ secara implisit}
jika dan hanya jika
\[
\Sigma(P) \cup \Sigma(P') \Entails \lforall[x_1][\dots
  \lforall[x_n][(\Atom{P}{x_1,\dots, x_n} \liff \Atom{P'}{x_1,\dots,
      x_n})]],
\]
dengan $\Sigma(P')$ sebagai hasil penggantian seragam $P$ oleh $P'$ dalam
$\Sigma(P)$.
\end{defn}

Dengan kata lain, bagi setiap model $\Struct{M}$ dan $R, R' \subseteq
\Domain{M}^n$, jika $\Expan{M}{R} \models \Sigma(P)$ dan
$\Expan{M}{R'} \models \Sigma(P')$, maka $R=R'$; di sini
$\Expan{M}{R}$ adalah !!{structure}~$\Struct{M'}$ bagi ekspansi
$\Lang{L}$ menjadi $\Lang{L} \cup \{P\}$ sedemikian sehingga
$\Assign{P}{M'} = R$, dan demikian pula bagi $\Expan{M}{R'}$.

\begin{thm}[Teorema Keterdefinisian Beth] Suatu himpunan $\Sigma(P)$ dari
  $\Lang{L} \cup\{P\}$-!!{sentence}s mendefinisikan $P$ secara implisit
  jika dan hanya jika $\Sigma(P)$ mendefinisikan $P$ secara eksplisit.
\end{thm}

\begin{proof}
Jika $\Sigma(P)$ mendefinisikan $P$ secara eksplisit, maka keduanya berlaku:
\begin{align*}
  \Sigma(P) & \Entails & \lforall[x_1][\dots \lforall[x_n]
    [(\Atom{P}{x_1,\dots, x_n} \liff !C(x_1,\dots,x_n))]]\\
  \Sigma(P') & \Entails & \lforall[x_1][\dots \lforall[x_n]
    [(\Atom{P'}{x_1,\dots, x_n} \liff !C(x_1,\dots,x_n))]]
\end{align*}
dan kesimpulannya mengikuti. Untuk arah sebaliknya, andaikan $\Sigma(P)$
mendefinisikan $P$ secara implisit. Mula-mula, kita menambahkan !!{constant}s
$c_1$, \dots,~$c_n$ ke $\Lang{L}$. Maka
\[
\Sigma(P) \cup \Sigma(P') \Entails
\Atom{P}{c_1, \dots, c_n} \to  \Atom{P'}{c_1, \dots, c_n}.
\]
Menurut kekompakan, terdapat himpunan hingga $\Delta_0 \subseteq \Sigma(P)$
dan $\Delta_1 \subseteq \Sigma(P')$ sedemikian sehingga
\[
\Delta_0 \cup \Delta_1 \Entails
\Atom{P}{c_1, \dots, c_n} \to \Atom{P'}{c_1, \dots, c_n}.
\]
Misalkan $!D(P)$ adalah konjungsi semua !!{sentence}s $!A(P)$ sedemikian
sehingga $!A(P) \in \Delta_0$ atau $!A(P') \in \Delta_1$, dan misalkan
$!D(P')$ adalah konjungsi semua !!{sentence}s $!A(P')$ sedemikian sehingga
$!A(P) \in \Delta_0$ atau $!A(P') \in \Delta_1$. Maka $!D(P) \land
!D(P') \Entails \Atom{P}{c_1, \dots, c_n} \to
\Atom{P'}{c_1,\dots,c_n}$. Kita dapat menyusun ulang ini sehingga setiap
!!{predicate} terdapat pada satu sisi $\Entails$:
\[
!D(P) \land \Atom{P}{c_1, \dots, c_n} \Entails
!D(P') \to \Atom{P'}{c_1, \dots, c_n}.
\]
Menurut Teorema Interpolasi Craig, terdapat !!a{sentence}
$!C(c_1,\dots, c_n)$ yang tidak memuat $P$ ataupun $P'$ sedemikian sehingga:
\begin{align*}
  !D(P) \land \Atom{P}{c_1, \dots, c_n} & \Entails !C(c_1,\dots, c_n); \\
  !C(c_1,\dots, c_n) & \Entails !D(P') \to \Atom{P'}{c_1, \dots, c_n}.
\end{align*}
Dari perikutan yang pertama, kita memperoleh $!D(P) \Entails
\Atom{P}{c_1,\dots, c_n} \lif !C(c_1,\dots, c_n)$. Dari yang kedua,
karena $\Expan{M}{R} \models !A(P)$ dalam bahasa $\Lang{L} \cup \{P\}$
jika dan hanya jika $\Expan{M}{R} \models !A(P')$ dalam bahasa
$\Lang{L} \cup \{P'\}$ yang bersesuaian, kita memperoleh
$!C(c_1,\dots, c_n) \Entails !D(P) \lif \Atom{P}{c_1,\dots, c_n}$,
dan dari sini:
\[
!D(P) \Entails !C(c_1,\dots,c_n) \to \Atom{P}{c_1,\dots, c_n}.
\]
Dengan menggabungkan keduanya, $!D(P) \Entails
\Atom{P}{c_1,\dots, c_n} \liff !C(c_1, \dots, c_n)$, dan menurut
monotonisitas serta generalisasi juga
\[
\Sigma(P) \Entails
\lforall[x_1][\dots\lforall[x_n][(\Atom{P}{x_1,\dots, x_n} \liff
    !C(x_1,\dots, x_n))]].
\]
\end{proof}

\end{document}
